\addtocontents{toc}{\protect\Needspace{3\baselineskip}}
\originalchapter{相对论运动与不可见的空间方向}\label{ch:basic-spacetime}
本章对应 8.251 第 1--2 讲. 公开笔记由 Bryan Owens 整理, 部分内容据 Kayla Jacobs 的课堂笔记; 教师为 Barton Zwiebach 与 Alan Guth. 以下重新组织说明, 并修正笔记中个别指标与根号的排印错误. 所需基础是向量内积、微分和一维薛定谔方程.

\originalsection{为什么用弦作为动力学对象}
普通绳子由原子组成, 它的张力和线密度可以独立测量. 强相互作用的通量管也可能在一定距离范围内表现为细长对象, 但其厚度和端点运动来自原来的场论. 基本弦理论则把一维延展对象作为出发点, 用不同振动模式产生不同粒子态. 这三种“弦”可共享面积作用量或波动方程, 却没有相同的微观假设.

我们先建立弦所处的目标时空, 再引入标记弦上的点的参数. 目标时空坐标的数目记为 $D$, 空间维数为 $d=D-1$. 本编前几章保留 $c,\hbar$, 进入量子谱后取 $c=\hbar=1$; 这比在同一公式中混用能量和逆长度更容易检查量纲.

\originalsection{间隔、指标与洛伦兹变换}
取 $x^0=ct$, 度规 $\eta_{\mu\nu}=\operatorname{diag}(-1,1,\ldots,1)$. 两事件的间隔是
\[
 \eta_{\mu\nu}\Delta x^\mu\Delta x^\nu
 =-c^2\Delta t^2+|\Delta\boldsymbol x|^2.
\]
它为负、零、正分别表示类时、类光、类空分离. 对类时轨迹另定义正的固有长度 $\dd s=\sqrt{-\eta_{\mu\nu}\dd x^\mu\dd x^\nu}$; 因而 $\dd s=c\dd t_{\rm prop}$. 有的笔记把“$\dd s^2$”直接叫洛伦兹间隔, 有的用它表示上述正长度平方, 必须先看定义.

降低指标是 $a_\mu=\eta_{\mu\nu}a^\nu$, 因此 $a_0=-a^0,a_i=a^i$, 内积 $a\cdot b=a_\mu b^\mu$. 下标不是新的向量分量列表, 而是同一向量通过度规得到的协向量. 沿 $x^1$ 的被动 boost 为
\[
 x^{\prime0}=\gamma(x^0-\beta x^1),\qquad
 x^{\prime1}=\gamma(x^1-\beta x^0),\qquad
 \gamma=(1-\beta^2)^{-1/2},\quad\beta=v/c.
\]
代回可见 $-(x^{\prime0})^2+(x^{\prime1})^2=-(x^0)^2+(x^1)^2$. 一般线性变换 $x'=\Lambda x$ 保持所有内积的充要条件为 $\Lambda^{\mathsf T}\eta\Lambda=\eta$. 这给出了洛伦兹群的定义, 也说明变换时不能只改变时间而保持同方向的空间坐标不变.

\originalsection{光锥时间的共轭量}
定义 $x^\pm=(x^0\pm x^1)/\sqrt2$, 横向方向用 $I=2,\ldots,D-1$ 表示. 则
\[
 \dd x\cdot\dd x=-2\dd x^+\dd x^-+\dd x^I\dd x^I,
 \quad\eta_{+-}=\eta_{-+}=-1,
 \quad a_+=-a^-,\quad a_-=-a^+.
\]
把 $x^+$ 用作演化参数不意味着与它共轭的能量就是 $p^+$. 平面波写成
\[
 \psi(x)=\exp(\ii p_\mu x^\mu/\hbar)
 =\exp\!\left[\frac{\ii}{\hbar}(-p^-x^+-p^+x^-+p^Ix^I)\right].
\]
故 $\ii\hbar\partial_+\psi=p^-\psi$. 此处 $x^+$ 的单位是长度, 所以其哈密顿量 $p^-$ 的单位是动量; 若把时间定义为 $x^+/c$, 对应能量为 $cp^-$. 普通时间对应 $E=cp^0$, 两者的角色来自相位中的共轭配对.

对质量 $m>0$ 的自由粒子, 固有时给 $u^\mu=\dd x^\mu/\dd t_{\rm prop}=\gamma(c,\boldsymbol v)$, $p^\mu=mu^\mu$. 于是
\[
 p\cdot p=-m^2c^2,\quad E^2=c^2|\boldsymbol p|^2+m^2c^4,
 \quad p^- =\frac{p^Ip^I+m^2c^2}{2p^+}.
\]
正能量类时粒子满足 $p^+>0$. 最右式让光锥哈密顿量成为横向动量的二次式, 是后续光锥量子化便利的原因. 类光且恰沿负 $x^1$ 运动的零模式可以有 $p^+=0$, 不能在该式中直接相除.

\originalsection{紧致空间不是可见的圆管表面}
在直线坐标上作识别 $y\sim y+2\pi R n$, $n\in\mathbb Z$, 得到周长 $2\pi R$ 的内禀圆. 这里没有被物理指定的圆心或外部径向方向. 基本域可取半开区间 $[0,2\pi R)$; 若画成闭区间, 两端代表同一点. 空間上的无限重复图像是覆盖空间中的同一个对象, 不是无限个独立物体.

两个独立识别 $(x,y)\sim(x+L_1,y)\sim(x,y+L_2)$ 给平坦矩形环面. 必须同时施加两种识别, 并在边界上补上相应配对. 基本域的面积 $L_1L_2$ 是环面的内禀面积.

\begin{center}
\begin{tikzpicture}[scale=.85,>=Stealth]
 \draw[->](-.3,0)--(6.6,0) node[right]{$y$};
 \draw[very thick,structurecolor](0,0)--(5,0);
 \fill (0,0) circle(2pt);\draw[fill=white](5,0) circle(2pt);
 \node[below]at(0,0){$0$};\node[below]at(5,0){$2\pi R$};
 \draw[->,bend left=35](0,.25)to node[above]{端点识别}(5,.25);
 \draw[xshift=8cm,->](0,0)--(2,0);
 \draw[xshift=8cm,->](0,0)--(1.41,1.41);
 \draw[xshift=8cm](1,0)arc(0:45:1);
 \node at(9.1,.48){$2\pi/n$};
 \node at(9,-.6){$n=8$ 的旋转商示意};
\end{tikzpicture}
\end{center}

\originalsection{商空间与固定点}
再识别 $y\sim-y$, 商 $\mathbb R/\mathbb Z_2$ 的基本域是半直线. 原点在反射下不动, 因而与一般点不同. 平面商 $z\sim\ee^{2\pi\ii/n}z$ 则有基本扇区 $0\leq\theta<2\pi/n$, 其两条射线应粘合. 局部平坦的圆锥具有亏角 $2\pi(1-1/n)$, 尖点处不能用覆盖平面上的普通坐标消除全局识别.

“紧致化”说明方向的长度有限; “轨形商”还保留群作用的固定点信息. 后续弦理论在固定点附近可能需要额外态, 所以几何识别并非只是一幅图的裁剪. 本章只建立基本域和动量量子化, 不把几何商自动当作已证明一致的量子弦背景.

\originalsection{低能实验为什么看不见小圆}
设非相对论粒子质量为 $m$, 沿 $0<x<a$ 有无限高势阱, 另有紧致方向 $y\sim y+2\pi R$. 对
\[
 H=-\frac{\hbar^2}{2m}(\partial_x^2+\partial_y^2)
\]
分离变量, 得归一化本征函数与能级
\[
 \psi_{k,n}=\sqrt{\frac{2}{a}}\sin\frac{k\pi x}{a}
 \frac{\ee^{\ii ny/R}}{\sqrt{2\pi R}},\qquad
 E_{k,n}=\frac{\hbar^2}{2m}\left(\frac{k^2\pi^2}{a^2}+\frac{n^2}{R^2}\right),
 \quad k=1,2,\ldots, n\in\mathbb Z.
\]
周期边界条件导致整数 $n$, 固定端点导致正整数 $k$. $n$ 与 $-n$ 可改写成正弦和余弦实基, 对 $n\ne0$ 为二重简并. 在 $k_{\rm B}T\ll\hbar^2/(2mR^2)$ 且实验碰撞能不足时, 激发紧致动量的代价很高, 主要只看到 $n=0$ 的一维谱. 这是有效降维, 空间本身并没有消失.

相对论粒子的相同周期性给 $p_y=n\hbar/R$, 因而在非紧致维度中
\[
 m_n^2=m_0^2+\frac{n^2\hbar^2}{c^2R^2}.
\]
这是 Kaluza--Klein 塔. 闭弦还允许绕圆的拓扑绕数, 其存在要从闭弦边界条件另行推导; 粒子的周期波函数本身没有绕弦能量.

\begin{example}
对电子质量的非相对论模型, 将圆半径减半使紧致激发能隙变为原来的四倍. 若能量已接近 $mc^2$, 上述薛定谔近似失效, 应用 $E^2=c^2p^2+m^2c^4$ 而非继续使用 $p^2/(2m)$.
\end{example}
